\documentclass[]{article}
% \documentclass[11pt]{article}

%%%%%%%%%%%%%%%%%%%%%
%% customized adds %%

\usepackage[round]{natbib}      % [author, year] --> (author, year)
\renewcommand\refname{}         % Change "Reference" section name --> None
\usepackage{amsmath,thmtools}
\usepackage{mathtools}          % sidebrace in equations
\usepackage{amssymb}            % for checkmark symbol, http://ctan.org/pkg/amssymb
\usepackage{tcolorbox}
\usepackage{multicol}
\usepackage[nottoc,numbib]{tocbibind} % Reference in TOC


%% --- page dimensions and margins (width - size .. etc) {
\usepackage[left=1.5in, right=1.0in, top=1.0in, bottom=1.0in]{geometry}
%% --- }


%% --- header-includes {
\usepackage{fancyhdr}
\pagestyle{fancy}
\fancyhead[CO,CE]{}
\fancyhead[]{}
\fancyfoot[L]{}
%% -- header --
% \chead{}
% \rhead{}
% \lhead{}
%% -- footers --
% \cfoot{} % keep to get rid of the bottom page
% \lfoot{\scriptsize PhD dissertation, \date{February 2019}}
% \rfoot{\scriptsize \author{XXXX XXXX}}
%% --- }



%% --- title page top-right logo {
% src: https://tex.stackexchange.com/a/184671

% \usepackage{eso-pic}
% \newcommand\AtPageUpperRight[1]{\AtPageUpperLeft{%
%    \makebox[\paperwidth][r]{#1}}}

%% --- }


%% --- Title page and committe approvals {
% \usepackage{titling}
% \pretitle{%
% }

% \posttitle{
% }
%% --- }


%% --- Remove dots and page number from toc {

%% https://tex.stackexchange.com/a/55466/81376
% \makeatletter
% \renewcommand{\@dotsep}{10000}
% \makeatother

%% https://tex.stackexchange.com/a/418828/81376
%\addtocontents{toc}{\protect\thispagestyle{empty}}
%\pagenumbering{gobble}

%% --- }



%% --- make each chapter title (i.e. section) starts in a new page {
\usepackage{titlesec}
\newcommand{\sectionbreak}{\clearpage}
%% --- }


%% --- double spacing lines {
% \usepackage{setspace}
% \doublespacing
%% --- }


%% --- indent a paragraph beginning {
\usepackage{setspace}\setlength\parindent{10pt}
%% --- }



%% --- algorithms env {

% \usepackage{amssymb}
\usepackage{float} % listing
\usepackage[linesnumbered,lined,boxed,commentsnumbered]{algorithm2e}
\floatstyle{ruled}
\newfloat{Listing}{thp}{lop}
\floatname{Listing}{Algorithm} % rename Listing -> Algorithm
%% setting keywords
\SetKwData{Left}{left}
\SetKwData{This}{this}
\SetKwData{Up}{up}
\SetKwFunction{Union}{Union}
\SetKwFunction{FindCompress}{FindCompress}
\SetKwFunction{uniform}{uniform}
\SetKwFunction{sample}{sample}
\SetKwFunction{context}{context}
\SetKwInOut{Input}{Input}
\SetKwInOut{Params}{Params}
\SetKwInOut{Output}{Output}

%% --- }


%% --- align quote rigth {
\usepackage{epigraph} % source: https://tex.stackexchange.com/a/65914/81376
\setlength{\epigraphwidth}{3.7in}
%% --- }

%%%%%%%%%%%%%

\usepackage{lmodern}
\usepackage{amssymb,amsmath}
\usepackage{ifxetex,ifluatex}
\usepackage{fixltx2e} % provides \textsubscript
\ifnum 0\ifxetex 1\fi\ifluatex 1\fi=0 % if pdftex
  \usepackage[T1]{fontenc}
  \usepackage[utf8]{inputenc}
\else % if luatex or xelatex
  \ifxetex
    \usepackage{mathspec}
  \else
    \usepackage{fontspec}
  \fi
  \defaultfontfeatures{Ligatures=TeX,Scale=MatchLowercase}
\fi
% use upquote if available, for straight quotes in verbatim environments
\IfFileExists{upquote.sty}{\usepackage{upquote}}{}
% use microtype if available
\IfFileExists{microtype.sty}{%
\usepackage{microtype}
\UseMicrotypeSet[protrusion]{basicmath} % disable protrusion for tt fonts
}{}
\usepackage[unicode=true,colorlinks]{hyperref} % color links
\hypersetup{
            pdftitle={Dissertation Title},
            pdfauthor={XXXX XXXX},
            pdfborder={0 0 0},
            breaklinks=true}
\urlstyle{same}  % don't use monospace font for urls
\usepackage{natbib}
\bibliographystyle{plainnat}
\usepackage{longtable,booktabs}
% Fix footnotes in tables (requires footnote package)
\IfFileExists{footnote.sty}{\usepackage{footnote}\makesavenoteenv{long table}}{}
\usepackage{graphicx,grffile}
\makeatletter

\def\maxwidth{\ifdim\Gin@nat@width>\linewidth\linewidth\else\Gin@nat@width\fi}
\def\maxheight{\ifdim\Gin@nat@height>\textheight\textheight\else\Gin@nat@height\fi}
\makeatother
% Scale images if necessary, so that they will not overflow the page
% margins by default, and it is still possible to overwrite the defaults
% using explicit options in \includegraphics[width, height, ...]{}
\setkeys{Gin}{width=\maxwidth,height=\maxheight,keepaspectratio}
\IfFileExists{parskip.sty}{%
\usepackage{parskip}
}{% else
\setlength{\parindent}{0pt}
\setlength{\parskip}{6pt plus 2pt minus 1pt}
}
\setlength{\emergencystretch}{3em}  % prevent overfull lines
\providecommand{\tightlist}{%
  \setlength{\itemsep}{0pt}\setlength{\parskip}{0pt}}
\setcounter{secnumdepth}{5}
% Redefines (sub)paragraphs to behave more like sections
\ifx\paragraph\undefined\else
\let\oldparagraph\paragraph
\renewcommand{\paragraph}[1]{\oldparagraph{#1}\mbox{}}
\fi
\ifx\subparagraph\undefined\else
\let\oldsubparagraph\subparagraph
\renewcommand{\subparagraph}[1]{\oldsubparagraph{#1}\mbox{}}
\fi

% set default figure placement to htbp
\makeatletter
\def\fps@figure{htbp}
\makeatother


\title{Dissertation Title}

\author{
by
\\~\\
XXXX XXXX
}

\providecommand{\institute}[1]{}
\institute{Pace University}
\date{February 2019}

\begin{document}
\renewcommand\headrule{} % remove fancyhdr line at the top
\pagenumbering{roman}    % starts roman numerials from the conver page
% Title page
\maketitle

\thispagestyle{empty} % remove styling (e.g. page number) from first page

%% --- document message {
\begin{center}
\vspace{3.5cm}  % line space
\textsc{
\large
Submitted in partial fulfillment \\
of the requirements for the degree of \\
Doctor of Philosophy in Computer Science
}
\\~\\
at
\\~\\
Seidenberg School of Computer Science and Information Systems
\\~\\
\textsc{Pace University}

\end{center}
%% --- }


%% --- Commmittee Signature page (approval sheet) {

\newpage
\thispagestyle{empty}


We hereby certify that this dissertation, submitted by \author{XXXX XXXX}, satisfies the dissertation requirements for the degree of Doctor of Philosophy in Computer Science and has been approved.

\vspace{2.5cm}  % line space

\newcommand{\signaturerule}{\rule{19em}{.4pt}}
\newcommand{\daterule}{\rule{8em}{.4pt}}
\renewcommand{\arraystretch}{2.0}

\begin{tabular}{lll}
    \bfseries Name  & \bfseries Signature & \bfseries Date \\
    \\
    \signaturerule &\daterule &\daterule\\
    Chair Person of Dissertation Committee\\
    \\
    \signaturerule &\daterule &\daterule\\
    Dissertation Committee\\
    \\
    \signaturerule &\daterule &\daterule\\
    Dissertation Committee\\

\end{tabular}

\vspace{4.5cm}  % line space

\begin{flushleft}
Seidenberg School of Computer Science and Information Systems
\\
Pace University, \date{February 2019}
\end{flushleft}


%% --- }


 % end title

\newpage
% \thispagestyle{empty}
{\setstretch{1.7} % abstract line spacing

\begin{abstract}
\large
A human being is a part of the whole called by us universe, a part
limited in time and space. He experiences himself, his thoughts and
feeling as something separated from the rest, a kind of optical delusion
of his consciousness. This delusion is a kind of prison for us,
restricting us to our personal desires and to affection for a few
persons nearest to us. Our task must be to free ourselves from this
prison by widening our circle of compassion to embrace all living
creatures and the whole of nature in its beauty. (By Einstein)
\end{abstract}
}
\newpage

% ---- Acknowledgements page {
\section*{\centering Acknowledgements}
\begin{center}
 % TODO: add your text below
Alhamdulillah
\end{center}
%% --- END of Acknowledgements text}

 % -- end abstract

{
% Table of Contents page TOC
\newpage
% \thispagestyle{empty}
\setcounter{tocdepth}{3}
\tableofcontents
\newpage
% \thispagestyle{empty}
\listoftables
\newpage
% \thispagestyle{empty}
\listoffigures
\newpage

}
 % end TOC

%% restore the regular header/footer and page numbering for the rest of the document
% \fancyhead[LO]{\leftmark}
\rhead{\thepage} % page number top right
\lhead[LO]{\leftmark} % section name top left
\cfoot{}
\lfoot{} %\scriptsize PhD dissertation, \date{February 2019}}
\rfoot{} %\scriptsize \author{XXXX XXXX}}
\pagenumbering{arabic} % restore page numering to Arabic
\setcounter{page}{1}
\doublespacing % make the document double spaced lines

% pandoc-xnos: macro to create a caption without a prefix
\makeatletter
\def\LT@makenoprefixcaption#1#2#3{%
  \LT@mcol\LT@cols c{\hbox to\z@{\hss\parbox[t]\LTcapwidth{
    \sbox\@tempboxa{#1{}#3}
    \ifdim\wd\@tempboxa>\hsize
      #1{}#3
    \else
      \hbox to\hsize{\hfil\box\@tempboxa\hfil}%
    \fi
    \endgraf\vskip\baselineskip}
  \hss}}}
\makeatother

% pandoc-tablenos: save original macros
\makeatletter
\let\LT@oldmakecaption=\LT@makecaption
\let\oldthetable=\thetable
\let\oldtheHtable=\theHtable
\makeatother

% pandoc-tablenos: environment disables table caption prefixes
\makeatletter
\newcounter{tableno}
\newenvironment{no-prefix-table-caption}{
  \let\LT@makecaption=\LT@makenoprefixcaption
  \renewcommand\thetable{x.\thetableno}
  \renewcommand\theHtable{x.\thetableno}
  \stepcounter{tableno}
}{
  \let\thetable=\oldthetable
  \let\theHtable=\oldtheHtable
  \let\LT@makecaption=\LT@oldmakecaption
  \addtocounter{table}{-1}
}
\makeatother

% pandoc-xnos: macro to create a caption without a prefix
\makeatletter
\long\def\@makenoprefixcaption#1#2{
  \vskip\abovecaptionskip
  \sbox\@tempboxa{#2}
  \ifdim \wd\@tempboxa >\hsize
    #2\par
  \else
    \global \@minipagefalse
    \hb@xt@\hsize{\hfil\box\@tempboxa\hfil}
  \fi
  \vskip\belowcaptionskip}
\makeatother

% pandoc-fignos: save original macros
\makeatletter
\let\@oldmakecaption=\@makecaption
\let\oldthefigure=\thefigure
\let\oldtheHfigure=\theHfigure
\makeatother

% pandoc-fignos: environment disables figure caption prefixes
\makeatletter
\newcounter{figno}
\newenvironment{no-prefix-figure-caption}{
  \let\@makecaption=\@makenoprefixcaption
  \renewcommand\thefigure{x.\thefigno}
  \renewcommand\theHfigure{x.\thefigno}
  \stepcounter{figno}
}{
  \let\thefigure=\oldthefigure
  \let\theHfigure=\oldtheHfigure
  \let\@makecaption=\@oldmakecaption
  \addtocounter{figure}{-1}
}
\makeatother

\section{Introduction}\label{introduction}

Some introduction goes here \ldots{}

\subsection{Informal introduction of the research problem with use cases
(why it is important, limitations of current solutions, your idea for
improving the solution
quality)}\label{informal-introduction-of-the-research-problem-with-use-cases-why-it-is-important-limitations-of-current-solutions-your-idea-for-improving-the-solution-quality}

Some informal introduction goes here \ldots{}

\subsection{Problem Statement (half page summary of research problem
with its scope and
objectives)}\label{problem-statement-half-page-summary-of-research-problem-with-its-scope-and-objectives}

Some problem statement goes here \ldots{}

\subsection{Solution methodology}\label{solution-methodology}

Some methodology description goes here \ldots{}

\subsection{Expected Contributions}\label{expected-contributions}

Some expected contribution goes here \ldots{}

\section{Literature Survey}\label{literature-survey}

\subsection{Minimum summary of all existing work supporting your
research}\label{minimum-summary-of-all-existing-work-supporting-your-research}

Some literature summary goes here \ldots{} Here is some cool reference
\citep{Goodfellow-et-al-2016}

\section{Problem Solution Design}\label{problem-solution-design}

Some solution goes here \ldots{}

\subsection{How to add a figure}\label{how-to-add-a-figure}

In markdown, this
\texttt{!{[}Pace\ University\ Logo{]}(assets/mylogo.png)} will generate
the following:



\begin{figure}
\centering
\includegraphics{assets/mylogo.png}
\caption{Pace University Logo}
\end{figure}



\subsection{How to add a table}\label{how-to-add-a-table}

In markdown, this

\begin{verbatim}
| a | b | c |
|:-:|:-:|:-:|
| 1 | 2 | 3 |
| 4 | 5 | 6 |
Table: Example table
\end{verbatim}

will generate the following table:



\begin{longtable}[]{@{}ccc@{}}
\caption{Example table}\tabularnewline
\toprule
a & b & c\tabularnewline
\midrule
\endfirsthead
\toprule
a & b & c\tabularnewline
\midrule
\endhead
1 & 2 & 3\tabularnewline
4 & 5 & 6\tabularnewline
\bottomrule
\end{longtable}



\subsection{How to cite}\label{how-to-cite}

It is possible to use Markdown/Pandoc style to cite any reference
e.g.~book, journal paper, conference \ldots{} etc. Just make sure to add
the reference's \texttt{BibTex} index to \texttt{assets/myref.bib} file,
then call its header as follows: \texttt{{[}@altowayan2016word{]}}.
LaTeX format is possible as well e.g.
\texttt{\textbackslash{}cite\{altowayan2016word\}} {[}\^{}footnote{]}

After compiling the final document it will generate this linked
reference \citep{altowayan2016word}

\subsection{How to change specific
format}\label{how-to-change-specific-format}

To change the layout or format of any component in the pdf document
e.g.~tables layout or figures position \ldots{} etc, one will need to
edit the LaTeX template file \texttt{assets/default.latex} to add and/or
modify the desired output.

\subsection{How to add a footnote}\label{how-to-add-a-footnote}

To add a footnote just use markdown flavored format. Here is an example:

Google\footnote{Google: \url{https://google.com}} is an awesome website
for learning just about anything.

In raw markdown, the above is written as follows:

\begin{verbatim}
Google[^mynote] is an awesome website for learning just about anything.
[^mynote]: Google: [https://google.com](https://google.com)
\end{verbatim}

\textbf{Fixed footers}:

To change the template's footers, look inside
\texttt{assets/default.latex} file for
\texttt{\textbackslash{}lfoot\{\}} and
\texttt{\textbackslash{}rfoot\{\}} commands. Add or change the needed
footer contents inside either command. For example, to make the date
appear at the bottom left:
\texttt{\textbackslash{}lfoot\{\textbackslash{}date\{\$date\$\}\}}

\section{Solution Continues}\label{solution-continues}

\subsection{Implementation Highlights (may include process to adapt your
solution for similar problems -- key for reusable
knowledge)}\label{implementation-highlights-may-include-process-to-adapt-your-solution-for-similar-problems-key-for-reusable-knowledge}

More solution goes here \ldots{}

\section{Experimental Validation}\label{experimental-validation}

Some validation goes here \ldots{}

\section{Conclusion}\label{conclusion}

Some conclusion goes here \ldots{}

\subsection{Contributions}\label{contributions}

Some contribution goes here \ldots{}

\subsection{Future Work}\label{future-work}

Some future work goes here \ldots{}

\renewcommand\refname{References}
\bibliography{/Users/Aziz/Dropbox/thesis/myref.bib}

\end{document}
